\section{Proofs}
\label{app:proofs}

\subsection{Empirical objective and population idealization}
For a training batch of $B$ windows and $H$ predicted time positions, the implemented loss is
\begin{equation}
\Lcal_{\mathrm{pred}}=\frac{1}{BHD}\sum_{b,h}\|\widehat z_{b,h+1}-z_{b,h+1}\|^2.
\end{equation}
The notation in Equation~\ref{eq:lewm} averages over time and uses this coordinate-mean convention.
One-step targets are teacher-forced at every temporal position, with a shorter history near the
start of a window. For encoded windows of $T$ frames, $P$ random unit projection directions $u_p$,
$K$ knots $\tau_k\in[0,3]$, and $w_k=e^{-\tau_k^2/2}$, the regularizer is
\begin{equation}
\label{eq:sigreg}
R=\frac{B}{TP}\sum_{t,p,k}q_k w_k
\left[\left(\frac1B\sum_b\cos(\tau_k u_p^\top z_{b,t})-w_k\right)^2
+\left(\frac1B\sum_b\sin(\tau_k u_p^\top z_{b,t})\right)^2\right],
\end{equation}
where $\Delta=3/(K-1)$, with $q_k=\Delta$ at the endpoints and $2\Delta$ in the interior.
The characteristic function is averaged across batch examples separately at each temporal position;
the factor $B$ is part of the implementation with $\lambda=0.09$.

The population statements use the nonnegative discrepancy
\begin{equation}
\label{eq:popsigreg}
\mathrm{SIGReg}_\infty(\mathcal L(z))=
\int_{\mathbb S^{D-1}}\!\int_0^3 w(\tau)
\left[\big(\E\cos(\tau\langle z,u\rangle)-w(\tau)\big)^2
+\big(\E\sin(\tau\langle z,u\rangle)\big)^2\right]d\tau\,d\sigma(u),
\end{equation}
where $w(\tau)=e^{-\tau^2/2}$ and $\sigma$ is uniform on the sphere.
Only the fact that this functional vanishes for a standard Gaussian is needed. It is not a theorem
about exact minima of the finite-batch objective.

\subsection{Nuisance optimum and the cost of innovations}
\paragraph{Proposition~\ref{prop:representation}.}
\begin{proof}
Let $\beta$ recover the nuisance from the observation, set $f^\star=\varphi\circ\beta$, and let
$g^\star$ copy the last latent. Then $z_t=\varphi(b)=z_{t+1}$, so every prediction residual vanishes.
The Gaussian marginal gives covariance $I_D$ and zero population regularizer. Both objective terms
are nonnegative, so their sum is globally minimized at zero. Training-distribution independence
gives $I(\varphi(b);p_{1:T})=0$.

For deployment, draw nuisance $B$ and the controller's random seed $U$ at reset, independently of
the initial state and transition randomness. Run the original controller in the environment with
nuisance $B$. Let a state-blind policy with private seed $(B,U)$ feed constant latent history
$\varphi(B)$ to a copy of that controller and reproduce its actions with $U$. Coupled to the same
initial state and transition randomness, the policies take the same first action. If their actions
agree through step $t$, their states, supplied latent histories, and seeds agree, so the next actions
agree as well. Induction gives identical trajectories, returns, and terminations. This does not
assume independence between $B$ and states subsequently influenced by its induced actions.
\end{proof}

\begin{corollary}[Innovation cost for uniformly state-resolving encoders]
\label{cor:strict}
Let $p_{t+1}=p_t+\mu(a_t)+w_t$, $\mathcal H_t=\sigma(p_{\le t},b,a_{\le t})$, and
$\Sigma_t=\Cov(w_t\mid\mathcal H_t)$. Assume finite second moments of states and encoded observations,
$x_t=R(p_t,b)$, fixed nuisance $b$, and
$\|f(R(p,b))-f(R(p',b))\|\ge\kappa\|p-p'\|$ for all admissible $p,p',b$, with $\kappa>0$.
Every predictor of latent and action history satisfies
\begin{equation}
\label{eq:innovation}
\frac1D\E\|g-f(x_{t+1})\|^2\ge\frac{\kappa^2}{D}\E\tr\Sigma_t.
\end{equation}
When $\E\tr\Sigma_t>0$, this exceeds the nuisance encoder's zero prediction loss under the
assumptions of Proposition~\ref{prop:representation}.
\end{corollary}
The assumption is stronger than injectivity or task sufficiency: every state distance must be
preserved at a uniform scale. Innovation is measured after conditioning on state history and nuisance;
persistent wind is not automatically fresh innovation.

\paragraph{Corollary~\ref{cor:strict}.}
\begin{proof}
Put $Z=f(x_{t+1})$ and $\mathcal C_t=\sigma(z_{\le t},a_{\le t})$.
Because observations and latents are measurable functions of state and nuisance, $\mathcal C_t\subseteq\mathcal H_t$.
Any predictor using a shorter window is also $\mathcal C_t$-measurable.
Conditional least squares and the tower property give
\begin{align*}
\E\|g-Z\|^2
&\ge\E\|Z-\E[Z\mid\mathcal C_t]\|^2\\
&=\E\|Z-\E[Z\mid\mathcal H_t]\|^2
 +\E\|\E[Z\mid\mathcal H_t]-\E[Z\mid\mathcal C_t]\|^2\\
&\ge\E\tr\Cov(Z\mid\mathcal H_t).
\end{align*}
The cross term vanishes because its first factor has conditional mean zero given $\mathcal H_t$.
Take conditionally independent copies $P,P'$ of $p_{t+1}$ given $\mathcal H_t$.
Such copies exist for Euclidean state spaces, using their regular conditional distributions.
The nuisance $b$ and current state $p_t$ are fixed under this conditioning.
The conditional independent-copy variance identity yields
\begin{align*}
\tr\Cov(Z\mid\mathcal H_t)
&=\tfrac12\E[\|f(R(P,b))-f(R(P',b))\|^2\mid\mathcal H_t]\\
&\ge\tfrac{\kappa^2}{2}\E[\|P-P'\|^2\mid\mathcal H_t]\\
&=\kappa^2\tr\Cov(p_{t+1}\mid\mathcal H_t)
=\kappa^2\tr\Sigma_t.
\end{align*}
The final equality follows because $p_t+\mu(a_t)$ is $\mathcal H_t$-measurable.
Taking expectations and dividing by $D$ proves Equation~\ref{eq:innovation}.
Averaging these inequalities gives the corresponding bound for a multi-time prediction objective.
Proposition~\ref{prop:representation} supplies the zero-loss comparator when its assumptions also hold.
\end{proof}

\subsection{Conditional prediction and geometric scale}
\begin{lemma}[Conditional mean]
\label{lem:mean}
For a fixed encoder and data law with $\E\|z'\|^2<\infty$, the unique squared-error minimizer
over measurable square-integrable predictors from history $c$ and action $a$ is
$m(c,a)=\E[z'\mid c,a]$, up to null sets.
\end{lemma}
\begin{proof}
Let $m=\E[z'\mid c,a]$.
For every square-integrable measurable $g(c,a)$, expand $g-z'=(g-m)+(m-z')$.
The cross term has expectation zero since $\E[m-z'\mid c,a]=0$.
Hence
\[
\E\|g-z'\|^2=\E\|g-m\|^2+\E\tr\Cov(z'\mid c,a).
\]
The second term does not depend on $g$.
The first is nonnegative and vanishes exactly when $g=m$ almost surely.
\end{proof}

\paragraph{The finite-sample regularizer.}
The finite-sample regularizer differs from Equation~\ref{eq:popsigreg}.
For independent identically distributed samples with characteristic function $\phi$ and empirical estimate $\widehat\phi_n$,
\begin{equation}
\E|\widehat\phi_n(t)-w(t)|^2
=|\phi(t)-w(t)|^2+\frac{1-|\phi(t)|^2}{n}.
\end{equation}
Indeed, $\E\widehat\phi_n=\phi$ and $\E|\widehat\phi_n-\phi|^2=(1-|\phi|^2)/n$, since each complex exponential has modulus one; expanding around $\phi$ gives the identity.
At a Gaussian law, the residual equals $(1-e^{-t^2})/n$.
This is estimator bias, not an optimization floor.
Changing the law changes both terms; dependent temporal samples introduce additional covariance terms. The implementation multiplies this statistic by batch size $B$, so its Gaussian finite-batch bias does not shrink as $1/B$ after that scaling.
Finite projections and knots also do not characterize Gaussianity.


\paragraph{Proposition~\ref{prop:amplitude}.}
\begin{proof}
Partition $\R$ into $I_j=[2\varepsilon j,2\varepsilon(j+1))$, $j\in\mathbb Z$.
Let $q_j=\Pr(G\in I_j)$ for $G\sim\Nrm(0,1)$.
Every $q_j$ is positive and their sum is one.
Let $F_j$ be the conditional Gaussian distribution function on $I_j$, with inverse quantile function $Q_j$ on $(0,1)$.
The conditional CDF is continuous and strictly increasing on the interval interior, so its inverse is measurable and takes values in that interior.

Draw $J$ with probabilities $q_j$, $U\sim\mathrm{Uniform}(0,1)$, and $V\sim\Nrm(0,I_{D-1})$, independently.
Set $b=(J,U,V)$; for $D=1$, omit $V$.
Take $R(p,b)=(p,b)$ and define
\begin{equation}
f_\varepsilon(R(p,b))=\big(Q_J((U+p)/2),V\big).
\end{equation}
The mixture $(U+p_t)/2$ is uniform on $(0,1)$, since $p_t$ is an independent fair bit.
Conditional on $J=j$, the first coordinate therefore has the Gaussian law restricted to $I_j$.
Mixing with weights $q_j$ recovers $\Nrm(0,1)$.
Independence of $V$ gives the joint law $\Nrm(0,I_D)$ and zero population regularizer.

Given the first coordinate $y$, its interval determines $J$ almost surely.
Write $v=F_J(y)$.
Then $p_t=\mathbf1\{v>1/2\}$ and $U=2v-p_t$ almost surely; the other coordinates recover $V$.
Thus the latent recovers both state and nuisance.
For fixed $b$, the two first coordinates are distinct and lie inside the same interval of width $2\varepsilon$.
All remaining coordinates agree, proving the distance bound.

Latent history reveals $b$, while the next bit remains independent and fair.
The next latent consequently has two equiprobable values, say $z^{(0)}(b)$ and $z^{(1)}(b)$.
Lemma~\ref{lem:mean} gives $m(b)=(z^{(0)}(b)+z^{(1)}(b))/2$.
Its conditional squared error is $\|z^{(1)}(b)-z^{(0)}(b)\|^2/4$, strictly positive and bounded by $\varepsilon^2$.
The midpoint is square-integrable by Jensen's inequality, since both state encodings have finite second moments.
Taking expectations and dividing by $D$ proves the coordinate-mean loss bound, including strict positivity since the conditional error is positive almost surely.
The midpoint depends only on $b$, so it is independent of $p_{t+1}$.
Finally, decoding the current bit permits $a_t=p_t$.
A state-blind action is independent of the current fair bit, giving success probability $1/2$.
\end{proof}


\section{The cube worlds}
\label{app:cubes}

The position-readout study crosses black, per-frame random, and per-episode random backgrounds
with deterministic or hidden-wind dynamics. Color-noise and black-background datasets contain
100,000 transitions; the gray-noise conditions contain 500,000. Only the four gray grid-JEPA
conditions have three independent model initializations. Table~\ref{tab:cube-main} and
Figure~\ref{fig:cube-readout} report that replicated comparison. Tables~\ref{tab:cube-readouts}
and~\ref{tab:cube-flow} retain seed-0 outcomes for the broader conditions, using the same canonical
benchmark files throughout. Differences across dataset sizes and model recipes are not isolated
nuisance effects.

\paragraph{Environment.}
The scene contains a red $6\times6$ object in a $63\times63$ RGB image with a red boundary. Four
cardinal actions move the object three pixels. Reset positions are uniformly sampled from integer
coordinates in $[16,42]^2$, and initial wind is zero. In wind conditions, each wind coordinate
independently increments by $-1$, $0$, or $1$ with probabilities $0.1$, $0.8$, and $0.1$, clipped
to $[-2,2]$; the updated wind is added to the action displacement. Background pixels are uniform
integers from 0 to 255; gray noise shares the same sample across RGB channels. Contact with the
boundary gives reward $-1$ and terminates; other rewards are zero, with a 128-step cap.

\paragraph{Read-out fitting.}
The nonlinear latent position probe is fitted on training data and evaluated on independent
\texttt{eval.npz} frames. Its spatial softmax-and-offset output is produced by two hidden layers of width 128, trained for
10,000 updates with batch 2,048 and Adam at initial learning rate $10^{-3}$ with cosine decay.
Grid inputs are standardized using training statistics. For class-token inputs, standardization and
PCA whitening are selected using the final 10\% of fitting rows, then the chosen variant is refitted
on the full fitting split. Whitening floors singular values at 1\% of the largest.
Applying this read-out
to predicted latents tests both prediction and read-out extrapolation. The pixel read-out used for
decoded flow images is a separate supervised CNN trained with synthetic background replacement
using ground-truth cube masks; it is a diagnostic, not a policy or a world-model training signal.
Rollout windows are sampled with replacement and overlap, and exclude death or reset during the
eight context and 16 forecast frames. The configured benchmark uses 2,048 windows for video metrics
and the first 512 for position diagnostics; exact sampled indices were not retained in the archived
result files. Surviving windows are a selected subset of wind trajectories, so these errors do not
establish death prediction or closed-loop control.



\begin{table}[t]
\centering
\caption{Seed-0 cube position L1 error (pixels) for the grid and class-token JEPAs. ``Real'' evaluates
held-out encoded frames; ``Pred. 8'' evaluates the same read-out on eight-step predicted latents.
Gray conditions use more training transitions than black and color conditions.}
\label{tab:cube-readouts}
\inputtable{cube_readouts.tex}
\end{table}

\begin{table}[t]
\centering
\caption{Seed-0 flow cube position L1 error (pixels). Latent and pixel columns use different
supervised read-outs. The persistent-color failure is visible before dynamics: real encoded position
has large error. These are model-recipe references rather than objective-only comparisons.}
\label{tab:cube-flow}
\inputtable{cube_flow.tex}
\end{table}

Earlier controller experiments use the three black/color background conditions. Their final
survival outcomes are retained below. The class-token JEPA's near-zero loss on per-episode texture
and poor reward-driven control are consistent with the known predictable-nuisance mechanism,
without exactly realizing Proposition~\ref{prop:representation}. The replicated gray-frame/wind
failure instead has high prediction loss and should not be interpreted as the same mechanism.
Still-world survival can be achieved by state-blind alternating actions, so it does not certify
state feedback.

\begin{table}[ht]
\centering
\caption{Survival (percent) in the supplementary cube worlds. ``LeWM'' denotes the class-token JEPA;
``grid'' denotes the grid JEPA; ``goal'' and ``heads'' use latent-goal and learned-return search,
respectively. PPO columns use the final saved policies; ``px'' and ``z'' identify pixel and latent
observations. Still-world survival does not certify state feedback.}
\label{tab:cube}
\inputtable{cube.tex}
\end{table}

\clearpage
\section{Hyperparameters and compute}
\label{app:hparams}

The gate and gray-noise cube conditions supply $500{,}000$ uniform-random transitions each. The earlier black and color-noise cube datasets contain $100{,}000$ transitions. Craftax uses its separate $500{,}000$-transition archive. Model training uses bfloat16 compute,
float32 parameters, and one A100. A $500$-step calibration estimates the training length fitting
$575$ seconds including compilation; the grid gate model trains for $8{,}500$ steps ($17{,}000$ and $34{,}000$ in the earlier budget extensions of Appendix~\ref{app:historical-gate}), the class-token cube models for $23{,}750$--$24{,}000$.
Model seeds $0$, $1$, $2$ of the grid JEPA differ only in the initialisation and data-order seed.

\paragraph{JEPA settings.}
Both JEPA architectures use batch $128$ and AdamW with learning rate and weight decay $10^{-3}$, $10\%$ warm-up
followed by cosine decay to zero, and gradient clipping at $1.0$. SIGReg uses weight $0.09$, $17$
knots on $[0,3]$, and $1024$ random projections. Neither architecture uses a stop-gradient, teacher, decoder, or
reward during world-model fitting. Linear--BatchNorm--GELU--linear projection heads flatten batch, time, and spatial tokens for normalization. Training-mode statistics therefore couple examples and temporal positions; at inference, running statistics define a fixed pointwise projection and the attention predictor is causal. The theoretical causal-map formulation is an inference/population idealization.

\begin{center}
{
\begin{tabular}{lll}
\toprule
& Class-token JEPA & Grid JEPA \\
\midrule
encoder            & ViT, $7\times7$ patches           & tile transformer $+$ conv \\
width / depth / heads & $192$ / $6$ / $3$              & $128$ / $6$ / $8$ \\
encoder MLP        & $768$                              & $4\times$ width \\
read-out           & class token                        & $4\times4$ grid, $32$ channels \\
latent dimension   & $192$                              & $512$ \\
projector          & BN projector, width $2048$      & same, hidden $512$ \\
predictor          & AdaLN-zero AR          & spatial + causal temporal \\
predictor depth / heads / dim & $6$ / $16$ / $64$  & $6$ / $8$ / $32$ \\
predictor MLP / width & $2048$ / latent                 & $2\times$ width / $256$ \\
history            & $3$ frames                         & $8$ frames \\
dropout            & $0.1$                              & $0$ \\
parameters         & $15.3$M                            & $2.6$M $+$ $8.3$M \\
\bottomrule
\end{tabular}}
\end{center}

\paragraph{Flow settings.}
The tokenizer uses $9{,}000$ steps, batch $256$, and learning rate $9\times10^{-4}$; dynamics
$16{,}000$ steps, batch $64$, the same learning rate, width $256$, six layers, eight spatial and
temporal attention heads, a $32$-frame context, independent per-frame noise levels, a shortcut loss of
weight $0.25$ every fourth step after the first quarter of training, history dropout $0.25$, two
history frames resampled every eight steps after half of training, and a temporal loss weight of
$0.25$. The symlog reward head has $255$ bins over $[-20,20]$; reward and termination losses have
weight $0.1$ and enter the trunk at full strength for half of training, then at factor $0.1$.
Inference uses four function evaluations per frame. We corrected terminal-transition masking in the
recipe's reward loss (Appendix~\ref{app:rewardbug}).

\paragraph{JEPA read-outs.}
Gate post-hoc heads use the frozen encoder and the first $800$ time rows ($80{,}000$ frames) of the same archive used for world-model training. The following $200$ time rows are held out from head fitting only; the independent world-model evaluation archive is \texttt{eval.npz}. The pixel decoder uses $8{,}000$ steps, batch $256$, learning
rate $10^{-3}$, width $256$, six residual blocks; it supports the figures and the decoded-latent control, in which its frozen output, clipped to $[0,255]$, replaces the pixels of the real-gate PPO; no true pixel enters that policy after the encoder. Reward,
termination, and behaviour-value heads are two-layer perceptrons of width $512$ on the latent and the
one-hot action, $6{,}000$ steps, batch $1024$, learning rate $10^{-3}$, Monte Carlo value targets at
discount $0.99$, value loss weight $0.1$. The residual predictor models $z_{t+1}-m$ with
$m=g(z_{t-7:t},a_{t-7:t})$ the original grid prediction, conditioned on $m$, $z_t$, and the action:
three residual blocks of width $1024$, $8{,}000$ steps, and batch $2048$. The final 20,000 valid
windows of the same world-model training archive are excluded from residual fitting and used for
its quality check. They are not independent world-model evaluation data. For encoder seed 0,
coordinate-mean MSE is $3.07\times10^{-4}$ for the original predictor, $2.05\times10^{-4}$ for one
residual sample, and $1.32\times10^{-4}$ for the mean of eight. This quality check uses eight function
evaluations per sample; the reported PPO runs instead use four. Averaging eight learned samples is
a random Monte Carlo average, not an exact conditional mean.

\paragraph{Policy optimisation.}
PPO uses discount $0.99$, GAE parameter $0.9$, clipping $0.2$, value coefficient $0.5$, entropy
coefficient $0.01$, AdamW at $2.5\times10^{-4}$ with weight decay $0.01$, four epochs per batch, and
gradient clipping at $1.0$. Pixels pass through a three-block Impala torso \citep{espeholt2018impala}, latents through a two-layer
perceptron; a GRU receives that embedding and the previous action; the critic uses $64$ categorical
bins with histogram-Gaussian targets \citep{farebrother2024stopregressing}. Real-environment PPO uses $128$ environments and $64$ steps per
iteration with four minibatch groups, truncated backpropagation over $16$ steps, and $1.5$M
environment steps. In-model PPO uses $150$ updates of $512$ trajectories of $16$ steps ($1.2$M
imagined steps) with minibatch size $128$; half of the starts are resets, half dataset states with
their real eight-frame context. All reported policy rows use final checkpoints.

\paragraph{Compute.} On a shared A100, the seed-0 grid-gate JEPA completes 8,500 updates in
573 seconds and has a logged total of 586.5 seconds. Probe and head fitting takes approximately
three minutes. The residual sampler's optimization stage takes 74.3 seconds; model loading, latent
extraction, fitting, and quality evaluation together take 275.2 seconds before the final file write.
Other logged stages are PPO inside the JEPA $2.0$--$4.3$ min per run; PPO inside the flow model
$5$--$6$ min per run at four function evaluations, $20$--$26$ min with the mean of eight; PPO in the real environment $3.7$--$7.7$ min on pixels with the convolutional torso and $51$--$70$ s with the MLP torso, $76$--$78$ s on the frozen latent, $4.7$--$4.8$ min on the decoded latent with the convolutional torso and $1.8$--$2.9$ min with the MLP torso; the latent diagnostic: $2$ min.

\section{Evaluation protocol}
\label{app:protocol}

Every gate-policy endpoint uses the final checkpoint and $256$ real episodes from reset, with the
$128$-step cap. Survival is the fraction still alive after step 128; death on that step is not
survival. Walls passed counts wall exits before or on the terminal transition, including an exit
on the same step as death. Discovery is final survival strictly above $50\%$.
Real-training and imagined-training runs advance their random keys under different iteration schedules; their final evaluations are not paired episode trials. Policy seeds are separate
runs conditional on a fitted model; independently fitted encoders provide a different replication
level. The flow model has only one checkpoint. Reported means weight policy runs equally;
standard deviations describe variation across those runs, not uncertainty across model families.

\paragraph{Gate specification.}
The gate uses the same image and object sizes and cardinal action displacement. A white vertical
wall of width three pixels spawns at $x=60$ every 22 steps and moves left three pixels per step.
Its 16-pixel gap is sampled at spawn, top or bottom with equal probability, and becomes visible
when the wall reaches $x\le40$. The gap location is hidden state before that reveal. Border or wall
contact produces reward $-1$ and termination; the time cap is 128 steps.

\paragraph{Latent diagnostic.} For each encoder, the diagnostic uses $64$ real contexts with the
wall still solid at $x\in\{42,45\}$. Reconstructed simulator states generate recorded futures and
matched wall-free futures under the recorded actions. The recorded branch survives the diagnostic
horizon by construction. At each step, wall amplitude is the squared Euclidean distance between
encodings of the wall-present and wall-absent frames. The teacher-forced residual uses real encoded
history and is a total squared error, summed over latent coordinates.
The wall coordinate of a prediction is its projection onto the wall-present minus wall-absent
latent difference, with endpoints $0$ and $1$. Decoded wall intensity averages the minimum colour
channel on shared wall pixels, excluding either possible gap band. These measurements compare
features within one encoder; raw distances across models are not an information metric.

\section{Latent-policy torso controls}
\label{app:latentctl}

The MLP on decoded latents passes $0.92$ walls per episode with no discoveries in six runs
(Table~\ref{tab:latentctl}), using encoders $0$ and $1$. The decoded CNN control instead uses all
three encoders, with three policies per encoder. Both of its discoveries use encoder $1$.
The decoder is frozen and reconstruction-supervised; success therefore establishes useful encoded
information under this added read-out, not sufficiency of the unsupervised training objective.
Failure of the true-pixel MLP recipe establishes no upper bound on MLP capacity.

\begin{table}[ht]
\centering
\caption{Real-gate observation and torso controls: final walls per episode and survival, mean
$\pm$ standard deviation over policies; discovery means final survival above $50\%$.
All policies use $1{,}499{,}136$ real transitions.}
\label{tab:latentctl}
\inputtable{latentctl.tex}
\end{table}

\section{Control outcomes and latent diagnostics}
\label{app:round2}

Table~\ref{tab:paired} gives aggregate final outcomes; Table~\ref{tab:seeds} lists every
imagined JEPA policy. The pooled imagined JEPA mean weights all original, sampled-residual and
mean-residual runs equally. The flow rows distinguish six sampled policies from three policies
using a mean of eight, all sharing the same model fit. They do not estimate variation across
independent flow-model fits. Tables~\ref{tab:seeddiag} and \ref{tab:latent} retain the encoder
and geometry diagnostics supporting the gate observations.

\begin{figure}[t]
\centering
\paperfig{control_by_seed.pdf}{1.0}
\caption{\textbf{Control outcomes by model fit and predictor variant.} Each point is a final policy
evaluated on 256 real episodes. Original prediction, residual sampling, and residual means are
separate conditions; model and policy replication are distinct. Horizontal bars show within-model means,
and the dashed line marks random actions. The flow reference uses one checkpoint.}
\label{fig:control}
\end{figure}

\begin{table}[ht]
\centering
\caption{Final control outcomes in the real gate, mean $\pm$ standard deviation over policy runs.
Imagined training uses $1{,}228{,}800$ transitions per policy; real training uses $1{,}499{,}136$.
Model-fit counts and policy-run counts identify distinct replication levels.}
\label{tab:paired}
\inputtable{paired.tex}
\end{table}

\begin{table}[ht]
\centering
\caption{Every imagined JEPA policy at its final checkpoint: survival (percent) and walls per
episode by encoder seed, predictor variant and PPO seed. Each cell uses $256$ real episodes.}
\label{tab:seeds}
\inputtable{seeds.tex}
\end{table}

\begin{table}[ht]
\centering
\caption{Grid-JEPA diagnostics by encoder seed. Wall amplitude and total teacher-forced residual
are squared distances at future steps $1$ and $8$; decoded intensity uses steps $1$ and $2$.
The wall coordinate has endpoints $0$ (wall absent) and $1$ (true encoding).
Head AUC compares each death cause with alive transitions; PPO wall counts pool predictor variants.}
\label{tab:seeddiag}
\inputtable{seeddiag.tex}
\end{table}

\begin{table}[ht]
\centering
\caption{JEPA geometry and decoded wall intensity for model seed $0$, averaged over $64$
pre-reveal contexts. Distances use matched actions and the same encoder; total residuals include
all latent coordinates. Wall intensity is measured on shared wall pixels.}
\label{tab:latent}
\inputtable{latent.tex}
\end{table}

\begin{figure}[t]
\centering
\paperfig{geometry_by_seed.pdf}{1.0}
\caption{\textbf{Within-encoder scale diagnostics.} For each of the three gate encoders, wall-removal
distance and teacher-forced total one-step residual are summed squared latent distances. Curves
average 64 pre-reveal contexts whose recorded futures survive the diagnostic horizon. These are
feature sensitivity and total prediction error, not commensurate measures of information content.}
\label{fig:geometry}
\end{figure}

\clearpage
\section{Terminal rewards in the flow recipe}
\label{app:rewardbug}

The original recipe masks transitions entering reset frames out of reward training. Here the only
nonzero reward occurs on terminal transitions, so that mask removes the reward signal.
The corrected head predicts the next reward from the current frame summary and action with
terminal transitions retained in its loss. The reported flow-policy outcomes use this corrected
convention. The single-checkpoint reference therefore denotes the complete tokenizer, dynamics,
reward and termination package, rather than an isolated change in transition stochasticity.


\section{Craftax protocols and single-checkpoint results}
\label{app:craftax}


\subsection{Fresh local-goal evaluation}
The new evaluation uses the same existing 512-dimensional checkpoint, model seed 0, with 128 starts
for each evaluation seed 101, 202, and 303. Goal generation and real execution have separate
pseudorandom streams; all controllers share initial states and the execution stream. The target is
the endpoint after six random actions, following seven warm-up actions from reset. The objective is
final player-position equality with survival throughout the six executed actions. Every evaluated
trajectory survives, so the survival requirement does not alter the measured rates. The target
position displacements are zero, one, two, or three grid cells in 137, 148, 84, and 15 cases,
respectively.

The open-loop planner solves once at horizon six. Receding-horizon planning receives each real
observation, executes the first selected action, and replans over the remaining horizons
$6,5,4,3,2,1$. Every CEM solve uses 512 candidate sequences, 32 elites, and six iterations. Thus
receding planning uses six times as many candidate-sequence searches and 3.5 times as many predicted
latent transitions as one-shot planning, excluding additional encoding. It is not a compute-matched
feedback intervention. Both controllers use terminal latent MSE and the native 17 actions, without
reward heads, decoder use, latent rescaling, or privileged coordinates in the planning cost.

\begin{table}[ht]
\centering
\caption{Craftax goal-success counts and Wilson 95\% intervals, conditional on one fitted model.
The reference replays the privileged sequence used to generate the goal under independent execution
noise; it is a protocol check, not an oracle or a guaranteed upper bound.}
\label{tab:craftax-counts}
\begin{tabular}{lrrr}
\toprule
Controller & Success / total & Rate (\%) & Wilson interval (\%) \\
\midrule
No-op & $137/384$ & $35.68$ & $[31.05,40.59]$ \\
Random & $65/384$ & $16.93$ & $[13.51,21.00]$ \\
JEPA open loop & $358/384$ & $93.23$ & $[90.26,95.34]$ \\
JEPA receding horizon & $369/384$ & $96.09$ & $[93.66,97.62]$ \\
Reference-action replay & $382/384$ & $99.48$ & $[98.12,99.86]$ \\
\bottomrule
\end{tabular}
\end{table}

Open-loop and receding-horizon success on displaced targets is $222/247$ and $232/247$; for
displacements of at least two cells it is $80/99$ and $87/99$. Random actions reach $19/247$ and
$1/99$ in these strata, while no-op reaches neither. Mean endpoint position L1 errors are $0.1068$
open loop, $0.0781$ receding horizon, $1.5599$ random, and $0.9401$ no-op, in grid-cell units.
Receding-horizon success exceeds open loop by 2.86 percentage points; a paired episode bootstrap
with 10,000 resamples gives $[0.78,4.95]$ percentage points. All intervals quantify evaluation
variation for this fixed checkpoint. The three evaluation seeds are not training replications.

\begin{table}[ht]
\centering
\caption{Craftax success by target displacement. The first column includes all 384 starts; the
others retain 247 displaced targets and 99 targets at least two grid cells away.}
\label{tab:craftax}
\inputtable{craftax.tex}
\end{table}

\subsection{Earlier exploratory goal and reward protocols}
The earlier results below use different evaluation episodes and a shared goal-generation/execution
noise protocol. They are not pooled with the fresh independent-noise evaluation above.

The grid checkpoint uses training seed 0, a tile encoder, $16\times32=512$ latent coordinates,
a six-layer factored predictor of width 256, history length eight, 17 actions, batch 128, and
8,500 updates. The recorded parameter count is 11.12 million. Training takes 604.5 seconds from
loading through the logged end-to-end timer (578.5 seconds in the update loop); thus it is an
approximately ten-minute run, not a strict sub-600-second result. The frozen decoder and heads add
175.6 seconds. Results are from one model fit and evaluation seed 0.

\paragraph{Goal generation and planning.}
Reset states receive seven random warm-up actions. Six further random actions generate a reachable
goal image. CEM optimizes a length-six action sequence in latent space, with population 512,
32 elites, and six iterations. The learned planner uses the goal image, without privileged position
coordinates or the generating actions in its cost. It executes the entire selected sequence without
observation feedback. Success is equality of final player positions. Goal generation, random actions,
and the selected plan use the same future simulator random keys; replaying the generating actions
therefore reaches all goals and serves only as a protocol check. This protocol samples nearby
reachable goals, not arbitrary distant tasks.

The planned endpoint position L1 error is $0.0781$, against $1.5313$ for random actions. Pixel MSE
on the 0--255 scale is $104.00$ versus $1477.92$. These quantities use the same 128 starts as the
reported $121/128$ versus $24/128$ reach counts. They do not quantify variation across fitted models.

\paragraph{Reward control and head generalization.}
Reward and value heads use the first 80,000 transitions (320 time rows across 250 streams) from the
world-model training archive. Held-out time rows are excluded from head fitting but not from encoder
training. Value fitting uses the first 80 time rows owing to truncated-return handling. Reward and
value fitting correlations are $0.922$ and $0.995$; held-out correlations are $0.353$ and $0.266$.
Held-out reward RMSE is $0.222$. Mean predicted held-out value is $2.356$, versus mean empirical
discounted return $1.348$.

Reward evaluation uses 128 starts, 100 steps, no autoreset, and rewards accumulated while alive.
The CEM controller executes one action and replans from each real observation. Random actions,
myopic fitted-reward argmax, one predicted step plus value, and CEM obtain returns
$1.4453$, $0.7813$, $0.9438$, and $0.4734$, respectively; their across-start standard deviations
are $1.0897$, $1.1109$, $1.2014$, and $1.1039$. Fractions alive at the end are $0.6328$, $0.7266$,
$0.6328$, and $0.6719$. Aggregate summaries do not preserve the per-start covariance needed for
paired confidence intervals. The archive's mean reward multiplied by 100 is not an online behavior
baseline and is excluded from the comparison.

Earlier Craftax experiments used a different $81\times32$ spatial latent and different horizons.
They are not merged with these 512-dimensional results. There is no matched online flow-controller
comparison for this checkpoint.
